\documentclass[11pt,a4paper]{article}
\usepackage[utf8]{inputenc}
\usepackage[T1]{fontenc}
\usepackage[spanish,es-minimal,es-noshorthands]{babel}
\AtBeginDocument{\renewcommand{\sin}{\operatorname{sen}}\renewcommand{\arcsin}{\operatorname{arcsen}}}
\usepackage{lmodern}
\usepackage[sfdefault,lining]{FiraSans}
\usepackage[a4paper,top=2.3cm,bottom=2.2cm,left=1.8cm,right=1.8cm,headheight=1.3cm,headsep=0.4cm,footskip=1cm]{geometry}
\usepackage{amsmath,amssymb}
\usepackage{icomma}
\usepackage{graphicx}
\usepackage[table]{xcolor}
\usepackage{tikz}
\usepackage{fontawesome5}
\usepackage{ccicons}
\usepackage{fancyhdr}
\usepackage{lastpage}
\usepackage{tasks}
\usepackage{needspace}
\usepackage[most]{tcolorbox}
\usepackage{array}
\usepackage[hidelinks]{hyperref}
\AtBeginDocument{\hypersetup{pdftitle={Ficha de ejercicios de la unidad 9: Estadística Unidimensional. 4º ESO},pdfsubject={Matemáticas A 4º ESO},pdfkeywords={estadística, población, muestra, muestreo, variables, tablas de frecuencias, datos agrupados, gráficos estadísticos, moda, media, mediana, rango, desviación media, varianza, desviación típica, coeficiente de variación, cuartiles, percentiles, diagrama de cajas}}}

\definecolor{navy}{HTML}{1B2748}
\definecolor{accent}{HTML}{E8603A}
\definecolor{teal}{HTML}{1F9E90}
\definecolor{soft}{HTML}{EEF1F8}
\definecolor{muted}{HTML}{5D6685}

\setlength{\parindent}{0pt}
\setlength{\parskip}{3pt}
\settasks{label=\alph*), label-format=\bfseries\color{navy}, label-width=1.3em, item-indent=1.8em, column-sep=0.6em, after-item-skip=5pt, before-skip=3pt, after-skip=4pt}

% ---------- capçalera i peu ----------
\newcommand{\llicencia}{\href{https://creativecommons.org/licenses/by-nc-sa/4.0/deed.es}{\ccbyncsa}\ \ Licencia Creative Commons Atribución-NoComercial-CompartirIgual 4.0 Internacional (CC BY-NC-SA 4.0)}
\pagestyle{fancy}
\fancyhf{}
\fancyhead[L]{\raisebox{-0.35\height}{\includegraphics[height=0.95cm]{logo-gva.png}}}
\fancyhead[C]{\small\color{muted}Ficha de ejercicios · Unidad 9: Estadística Unidimensional · 4\textsuperscript{o} ESO}
\fancyhead[R]{\raisebox{-0.35\height}{\includegraphics[height=0.95cm]{logo-vermellar.png}}}
\fancyfoot[L]{\scriptsize\color{muted}\llicencia}
\fancyfoot[R]{\small\color{muted}\thepage}
\renewcommand{\headrulewidth}{0.4pt}
\renewcommand{\headrule}{\color{muted!40}\hrule height \headrulewidth}
\fancypagestyle{primera}{\fancyhf{}\renewcommand{\headrulewidth}{0pt}
  \fancyfoot[L]{\scriptsize\color{muted}\llicencia}
  \fancyfoot[R]{\small\color{muted}\thepage}}

% ---------- nivells i exercicis ----------
\newcommand{\estrelles}[1]{{\color{accent}\ifcase#1\or\faStar\faStar[regular]\faStar[regular]\or\faStar\faStar\faStar[regular]\or\faStar\faStar\faStar\fi}}
\newcounter{exer}
\newcommand{\ex}[1]{\par\needspace{4\baselineskip}\medskip\stepcounter{exer}\noindent{\color{navy}\bfseries\theexer.}\ \ {\small\estrelles{#1}}\ \ }

\newcommand{\seccio}[2]{\par\needspace{10\baselineskip}\bigskip
  \begin{tcolorbox}[colback=navy,colframe=navy,coltext=white,arc=2mm,boxrule=0pt,left=3mm,right=3mm,top=1.2mm,bottom=1.2mm]
  {\large\bfseries #1\quad #2}
  \end{tcolorbox}\vspace{-2pt}}

\newcolumntype{C}[1]{>{\centering\arraybackslash}p{#1}}
\definecolor{violet}{HTML}{7C6CF0}
\definecolor{gold}{HTML}{E0A21B}
% quadre de recordatori (només a la fitxa d'exercicis)
\newcommand{\recorda}[1]{\par\smallskip\begin{tcolorbox}[colback=accent!7,colframe=accent!55,arc=2mm,boxrule=0.5pt,left=2.5mm,right=2.5mm,top=1mm,bottom=1mm]\small{\color{accent}\faLightbulb[regular]}\ \textbf{Recuerda:}\ #1\end{tcolorbox}}
\newcommand{\res}[1]{#1}\newcommand{\rest}[1]{#1}

\def\fmneed{14}
% =====================================================================
% Capçalera FentMates de les fitxes: un color per a cada curs.
% Els colors són els de la pàgina inicial (LEVELS a assets/js/materials.js);
% el logo ve de _eines/logo/pdf/ (python3 _eines/logo/genera.py).
%
% Ús, al final del preàmbul de la unitat:
%   % =====================================================================
% Capçalera FentMates de les fitxes: un color per a cada curs.
% Els colors són els de la pàgina inicial (LEVELS a assets/js/materials.js);
% el logo ve de _eines/logo/pdf/ (python3 _eines/logo/genera.py).
%
% Ús, al final del preàmbul de la unitat:
%   % =====================================================================
% Capçalera FentMates de les fitxes: un color per a cada curs.
% Els colors són els de la pàgina inicial (LEVELS a assets/js/materials.js);
% el logo ve de _eines/logo/pdf/ (python3 _eines/logo/genera.py).
%
% Ús, al final del preàmbul de la unitat:
%   \input{../comu/fentmates.tex}
%   \fmcurs{1eso}
% i, al començament del document:
%   \capcalera{Matemàtiques}{1r ESO}{Fitxa d'exercicis · Unitat 2}{Divisibilitat}
% Opcional abans de \input: \def\fmneed{14} (espai mínim abans de cada secció).
% =====================================================================
\makeatletter
\def\fm@logodir{../../logo/pdf}
\providecommand{\fmneed}{10}
\def\fm@cursnom{}

% clau del curs, color i color del text damunt del color
\newcommand{\fm@curs}[3]{\@namedef{fm@color@#1}{#2}\@namedef{fm@text@#1}{#3}}
\fm@curs{1eso}{E6B84A}{navy}
\fm@curs{2eso}{3B82F6}{white}
\fm@curs{3eso}{FF7A50}{white}
\fm@curs{4eso-a}{A78BFA}{white}
\fm@curs{4eso-b}{6D28D9}{white}
\fm@curs{1batx-mi}{9B1C1C}{white}
\fm@curs{1batx-ccss}{EF6B6B}{white}
\fm@curs{2batx-mii}{1E6B47}{white}
\fm@curs{2batx-ccss}{2BB3A3}{white}

\newcommand{\fmcurs}[1]{%
  \@ifundefined{fm@color@#1}{\PackageError{fentmates}{Curs desconegut: #1}{}}{}%
  \edef\fm@tmp{\noexpand\definecolor{curs}{HTML}{\@nameuse{fm@color@#1}}%
               \noexpand\colorlet{curstext}{\@nameuse{fm@text@#1}}}\fm@tmp
  \colorlet{cursfosc}{curs!72!black}%
  \def\fm@logo{\fm@logodir/vertical-#1.pdf}%
  \def\fm@marca{\fm@logodir/marca-#1.pdf}}

% ---------- capçalera de la primera pàgina ----------
% #1 matèria (petit), #2 curs, #3 tipus de document i unitat, #4 títol
\newcommand{\capcalera}[4]{%
  \gdef\fm@cursnom{#2}%
  \par\noindent
  \begin{tcolorbox}[colback=curs,colframe=curs,boxrule=0pt,arc=2.2mm,width=3.5cm,height=2.15cm,
      valign=center,left=3mm,right=2mm,top=1mm,bottom=1mm,nobeforeafter,box align=center]
    \color{curstext}{\footnotesize #1}\par\vspace{1pt}{\fontsize{19}{21}\selectfont\bfseries #2\par}%
  \end{tcolorbox}\hspace{5mm}%
  \begin{minipage}[c]{\dimexpr\linewidth-3.5cm-5mm-2.9cm\relax}
    {\color{cursfosc}\small\bfseries\MakeUppercase{#3}}\par\vspace{2pt}
    {\color{navy}\fontsize{20}{23}\selectfont\bfseries #4\par}
  \end{minipage}\hfill
  \begin{minipage}[c]{2.6cm}\raggedleft\includegraphics[width=2.6cm]{\fm@logo}\end{minipage}\par
  \vspace{2pt}}

% ---------- barres de secció del color del curs ----------
\renewcommand{\seccio}[2]{\par\needspace{\fmneed\baselineskip}\bigskip
  \begin{tcolorbox}[colback=curs,colframe=curs,coltext=curstext,arc=2mm,boxrule=0pt,left=3mm,right=3mm,top=1.2mm,bottom=1.2mm]
  {\large\bfseries #1\quad #2}
  \end{tcolorbox}\vspace{-2pt}}

% ---------- estrelles de dificultat del color del curs (a Batxillerat no se'n posen) ----------
\@ifundefined{estrelles}{}{%
  \renewcommand{\estrelles}[1]{{\color{curs}\ifcase#1\or\faStar\faStar[regular]\faStar[regular]\or\faStar\faStar\faStar[regular]\or\faStar\faStar\faStar\fi}}}

% ---------- capçalera i peu de les pàgines ----------
% primera pàgina: logos institucionals al peu
\fancypagestyle{primera}{\fancyhf{}\renewcommand{\headrulewidth}{0pt}
  \fancyfoot[L]{\raisebox{-0.32\height}{\includegraphics[height=0.78cm]{logo-gva.png}}\hspace{5mm}%
    \raisebox{-0.32\height}{\includegraphics[height=0.78cm]{logo-vermellar.png}}}
  \fancyfoot[R]{\scriptsize\color{muted}\href{https://creativecommons.org/licenses/by-nc-sa/4.0/deed.ca}{\ccbyncsa}\ \ CC BY-NC-SA 4.0 · fentmates.org\qquad\small\thepage}}
% la resta: el símbol a l'esquerra i el curs a la dreta, amb el títol al centre com abans
\fancyhead[L]{\raisebox{-0.3\height}{\includegraphics[height=0.62cm]{\fm@marca}}}
\fancyhead[R]{\small\bfseries\color{cursfosc}\fm@cursnom}
\renewcommand{\headrule}{\color{curs!45}\hrule height \headrulewidth}
\makeatother

%   \fmcurs{1eso}
% i, al començament del document:
%   \capcalera{Matemàtiques}{1r ESO}{Fitxa d'exercicis · Unitat 2}{Divisibilitat}
% Opcional abans de \input: \def\fmneed{14} (espai mínim abans de cada secció).
% =====================================================================
\makeatletter
\def\fm@logodir{../../logo/pdf}
\providecommand{\fmneed}{10}
\def\fm@cursnom{}

% clau del curs, color i color del text damunt del color
\newcommand{\fm@curs}[3]{\@namedef{fm@color@#1}{#2}\@namedef{fm@text@#1}{#3}}
\fm@curs{1eso}{E6B84A}{navy}
\fm@curs{2eso}{3B82F6}{white}
\fm@curs{3eso}{FF7A50}{white}
\fm@curs{4eso-a}{A78BFA}{white}
\fm@curs{4eso-b}{6D28D9}{white}
\fm@curs{1batx-mi}{9B1C1C}{white}
\fm@curs{1batx-ccss}{EF6B6B}{white}
\fm@curs{2batx-mii}{1E6B47}{white}
\fm@curs{2batx-ccss}{2BB3A3}{white}

\newcommand{\fmcurs}[1]{%
  \@ifundefined{fm@color@#1}{\PackageError{fentmates}{Curs desconegut: #1}{}}{}%
  \edef\fm@tmp{\noexpand\definecolor{curs}{HTML}{\@nameuse{fm@color@#1}}%
               \noexpand\colorlet{curstext}{\@nameuse{fm@text@#1}}}\fm@tmp
  \colorlet{cursfosc}{curs!72!black}%
  \def\fm@logo{\fm@logodir/vertical-#1.pdf}%
  \def\fm@marca{\fm@logodir/marca-#1.pdf}}

% ---------- capçalera de la primera pàgina ----------
% #1 matèria (petit), #2 curs, #3 tipus de document i unitat, #4 títol
\newcommand{\capcalera}[4]{%
  \gdef\fm@cursnom{#2}%
  \par\noindent
  \begin{tcolorbox}[colback=curs,colframe=curs,boxrule=0pt,arc=2.2mm,width=3.5cm,height=2.15cm,
      valign=center,left=3mm,right=2mm,top=1mm,bottom=1mm,nobeforeafter,box align=center]
    \color{curstext}{\footnotesize #1}\par\vspace{1pt}{\fontsize{19}{21}\selectfont\bfseries #2\par}%
  \end{tcolorbox}\hspace{5mm}%
  \begin{minipage}[c]{\dimexpr\linewidth-3.5cm-5mm-2.9cm\relax}
    {\color{cursfosc}\small\bfseries\MakeUppercase{#3}}\par\vspace{2pt}
    {\color{navy}\fontsize{20}{23}\selectfont\bfseries #4\par}
  \end{minipage}\hfill
  \begin{minipage}[c]{2.6cm}\raggedleft\includegraphics[width=2.6cm]{\fm@logo}\end{minipage}\par
  \vspace{2pt}}

% ---------- barres de secció del color del curs ----------
\renewcommand{\seccio}[2]{\par\needspace{\fmneed\baselineskip}\bigskip
  \begin{tcolorbox}[colback=curs,colframe=curs,coltext=curstext,arc=2mm,boxrule=0pt,left=3mm,right=3mm,top=1.2mm,bottom=1.2mm]
  {\large\bfseries #1\quad #2}
  \end{tcolorbox}\vspace{-2pt}}

% ---------- estrelles de dificultat del color del curs (a Batxillerat no se'n posen) ----------
\@ifundefined{estrelles}{}{%
  \renewcommand{\estrelles}[1]{{\color{curs}\ifcase#1\or\faStar\faStar[regular]\faStar[regular]\or\faStar\faStar\faStar[regular]\or\faStar\faStar\faStar\fi}}}

% ---------- capçalera i peu de les pàgines ----------
% primera pàgina: logos institucionals al peu
\fancypagestyle{primera}{\fancyhf{}\renewcommand{\headrulewidth}{0pt}
  \fancyfoot[L]{\raisebox{-0.32\height}{\includegraphics[height=0.78cm]{logo-gva.png}}\hspace{5mm}%
    \raisebox{-0.32\height}{\includegraphics[height=0.78cm]{logo-vermellar.png}}}
  \fancyfoot[R]{\scriptsize\color{muted}\href{https://creativecommons.org/licenses/by-nc-sa/4.0/deed.ca}{\ccbyncsa}\ \ CC BY-NC-SA 4.0 · fentmates.org\qquad\small\thepage}}
% la resta: el símbol a l'esquerra i el curs a la dreta, amb el títol al centre com abans
\fancyhead[L]{\raisebox{-0.3\height}{\includegraphics[height=0.62cm]{\fm@marca}}}
\fancyhead[R]{\small\bfseries\color{cursfosc}\fm@cursnom}
\renewcommand{\headrule}{\color{curs!45}\hrule height \headrulewidth}
\makeatother

%   \fmcurs{1eso}
% i, al començament del document:
%   \capcalera{Matemàtiques}{1r ESO}{Fitxa d'exercicis · Unitat 2}{Divisibilitat}
% Opcional abans de \input: \def\fmneed{14} (espai mínim abans de cada secció).
% =====================================================================
\makeatletter
\def\fm@logodir{../../logo/pdf}
\providecommand{\fmneed}{10}
\def\fm@cursnom{}

% clau del curs, color i color del text damunt del color
\newcommand{\fm@curs}[3]{\@namedef{fm@color@#1}{#2}\@namedef{fm@text@#1}{#3}}
\fm@curs{1eso}{E6B84A}{navy}
\fm@curs{2eso}{3B82F6}{white}
\fm@curs{3eso}{FF7A50}{white}
\fm@curs{4eso-a}{A78BFA}{white}
\fm@curs{4eso-b}{6D28D9}{white}
\fm@curs{1batx-mi}{9B1C1C}{white}
\fm@curs{1batx-ccss}{EF6B6B}{white}
\fm@curs{2batx-mii}{1E6B47}{white}
\fm@curs{2batx-ccss}{2BB3A3}{white}

\newcommand{\fmcurs}[1]{%
  \@ifundefined{fm@color@#1}{\PackageError{fentmates}{Curs desconegut: #1}{}}{}%
  \edef\fm@tmp{\noexpand\definecolor{curs}{HTML}{\@nameuse{fm@color@#1}}%
               \noexpand\colorlet{curstext}{\@nameuse{fm@text@#1}}}\fm@tmp
  \colorlet{cursfosc}{curs!72!black}%
  \def\fm@logo{\fm@logodir/vertical-#1.pdf}%
  \def\fm@marca{\fm@logodir/marca-#1.pdf}}

% ---------- capçalera de la primera pàgina ----------
% #1 matèria (petit), #2 curs, #3 tipus de document i unitat, #4 títol
\newcommand{\capcalera}[4]{%
  \gdef\fm@cursnom{#2}%
  \par\noindent
  \begin{tcolorbox}[colback=curs,colframe=curs,boxrule=0pt,arc=2.2mm,width=3.5cm,height=2.15cm,
      valign=center,left=3mm,right=2mm,top=1mm,bottom=1mm,nobeforeafter,box align=center]
    \color{curstext}{\footnotesize #1}\par\vspace{1pt}{\fontsize{19}{21}\selectfont\bfseries #2\par}%
  \end{tcolorbox}\hspace{5mm}%
  \begin{minipage}[c]{\dimexpr\linewidth-3.5cm-5mm-2.9cm\relax}
    {\color{cursfosc}\small\bfseries\MakeUppercase{#3}}\par\vspace{2pt}
    {\color{navy}\fontsize{20}{23}\selectfont\bfseries #4\par}
  \end{minipage}\hfill
  \begin{minipage}[c]{2.6cm}\raggedleft\includegraphics[width=2.6cm]{\fm@logo}\end{minipage}\par
  \vspace{2pt}}

% ---------- barres de secció del color del curs ----------
\renewcommand{\seccio}[2]{\par\needspace{\fmneed\baselineskip}\bigskip
  \begin{tcolorbox}[colback=curs,colframe=curs,coltext=curstext,arc=2mm,boxrule=0pt,left=3mm,right=3mm,top=1.2mm,bottom=1.2mm]
  {\large\bfseries #1\quad #2}
  \end{tcolorbox}\vspace{-2pt}}

% ---------- estrelles de dificultat del color del curs (a Batxillerat no se'n posen) ----------
\@ifundefined{estrelles}{}{%
  \renewcommand{\estrelles}[1]{{\color{curs}\ifcase#1\or\faStar\faStar[regular]\faStar[regular]\or\faStar\faStar\faStar[regular]\or\faStar\faStar\faStar\fi}}}

% ---------- capçalera i peu de les pàgines ----------
% primera pàgina: logos institucionals al peu
\fancypagestyle{primera}{\fancyhf{}\renewcommand{\headrulewidth}{0pt}
  \fancyfoot[L]{\raisebox{-0.32\height}{\includegraphics[height=0.78cm]{logo-gva.png}}\hspace{5mm}%
    \raisebox{-0.32\height}{\includegraphics[height=0.78cm]{logo-vermellar.png}}}
  \fancyfoot[R]{\scriptsize\color{muted}\href{https://creativecommons.org/licenses/by-nc-sa/4.0/deed.ca}{\ccbyncsa}\ \ CC BY-NC-SA 4.0 · fentmates.org\qquad\small\thepage}}
% la resta: el símbol a l'esquerra i el curs a la dreta, amb el títol al centre com abans
\fancyhead[L]{\raisebox{-0.3\height}{\includegraphics[height=0.62cm]{\fm@marca}}}
\fancyhead[R]{\small\bfseries\color{cursfosc}\fm@cursnom}
\renewcommand{\headrule}{\color{curs!45}\hrule height \headrulewidth}
\makeatother

\fmcurs{4eso-a}

\begin{document}
\thispagestyle{primera}

\capcalera{Matemáticas A}{4\textsuperscript{o} ESO}{Ficha de ejercicios · Unidad 9}{Estadística Unidimensional}
\vspace{10pt}
Nombre y apellidos: \rule{6.6cm}{0.4pt}\hfill Grupo: \rule{1.6cm}{0.4pt}\hfill Fecha: \rule{2.2cm}{0.4pt}
\vspace{8pt}

\begin{tcolorbox}[colback=soft,colframe=soft,arc=2mm,left=3mm,right=3mm,top=1.5mm,bottom=1.5mm]
\small Ejercicios graduados por niveles de dificultad:\quad \estrelles{1}\ básico \qquad \estrelles{2}\ medio \qquad \estrelles{3}\ avanzado.\\[2pt] Resuélvelos en el cuaderno mostrando todos los pasos. Puedes usar la calculadora en modo estadístico para comprobar los resultados. Las secciones siguen el orden de los apuntes de la unidad (9.1--9.8).
\end{tcolorbox}

% =====================================================
\seccio{1}{Población, muestra y variables estadísticas}

\ex{1} Clasifica estas variables en cualitativas, cuantitativas discretas o cuantitativas continuas:\begin{tasks}(2)
  \task número de mensajes enviados en un día
  \task peso de una maleta
  \task marca del coche familiar
  \task tiempo de espera en la parada del autobús
  \task número de hermanos
  \task nivel de estudios (ESO, Bachillerato…)
\end{tasks}\recorda{Si los valores no son números, es \textbf{cualitativa}. Si son números que se \textbf{cuentan}, es discreta; si se \textbf{miden}, es continua.}

\ex{1} Para saber cuántas horas de deporte hacen a la semana los 900 alumnos de un instituto, se encuesta a 90 elegidos al azar. Indica la población, la muestra y su tamaño, un individuo y la variable estudiada (y de qué tipo es).

\ex{2} Razona si estas muestras son representativas:\begin{tasks}(1)
  \task Para conocer el deporte preferido de los jóvenes de una ciudad, se pregunta a 50 personas que salen de un campo de fútbol.
  \task Para conocer la altura media del alumnado de un instituto, se eligen al azar 5 alumnos de cada grupo.
\end{tasks}

\ex{2} Un pueblo tiene 2\,400 habitantes: 600 menores de 18 años, 1\,200 de entre 18 y 65 años y 600 de más de 65. Se quiere hacer una encuesta a 120 personas con un muestreo estratificado. ¿Cuántas personas hay que elegir de cada grupo de edad?\recorda{En el muestreo \textbf{estratificado}, cada grupo aporta a la muestra la misma proporción que tiene en la población.}

\ex{3} De la lista de los 840 alumnos de un instituto queremos elegir 60 con un muestreo sistemático. ¿Cada cuántos alumnos tenemos que coger uno? Si el primero elegido al azar es el número 9, ¿cuáles son el segundo, el quinto y el último?


% =====================================================
\seccio{2}{Tablas de frecuencias con datos aislados}

\ex{1} Estos son los libros que han leído en un mes 24 alumnos: \ $1,\ \allowbreak 2,\ \allowbreak 4,\ \allowbreak 2,\ \allowbreak 0,\ \allowbreak 1,\ \allowbreak 3,\ \allowbreak 2,\ \allowbreak 2,\ \allowbreak 2,\ \allowbreak 5,\ \allowbreak 2,\ \allowbreak 1,\ \allowbreak 1,\ \allowbreak 3,\ \allowbreak 4,\ \allowbreak 0,\ \allowbreak 0,\ \allowbreak 2,\ \allowbreak 2,\ \allowbreak 1,\ \allowbreak 3,\ \allowbreak 3,\ \allowbreak 1$. \ Haz la tabla de frecuencias con $f_i$, $F_i$, $h_i$ y porcentajes.

\ex{1} Con la tabla del ejercicio anterior, responde:\begin{tasks}(2)
  \task ¿Cuántos alumnos han leído como mucho 2 libros?
  \task ¿Cuántos han leído más de 3?
  \task ¿Qué porcentaje ha leído 3 o menos?
  \task ¿Qué porcentaje no ha leído ninguno?
\end{tasks}

\ex{2} Completa la tabla del deporte preferido de 40 alumnos:\par\smallskip\begin{center}{\renewcommand{\arraystretch}{1.25}\begin{tabular}{|c|C{1.2cm}|C{1.2cm}|C{1.2cm}|}\hline Deporte & $f_i$ & $h_i$ & \%\\\hline Fútbol & 14 & & \\\hline Baloncesto & & $0{,}25$ & \\\hline Tenis & & & 15\\\hline Natación & 8 & & \\\hline Otros & & & \\\hline \textbf{Total} & 40 & & \\\hline\end{tabular}}\end{center}

\ex{2} En esta tabla, $x_i$ es el número de visitas al dentista en un año y $F_i$, la frecuencia acumulada. Encuentra las frecuencias absolutas. ¿Cuántas personas han ido exactamente 3 veces? ¿Y al menos 2 veces?\par\smallskip\begin{center}{\renewcommand{\arraystretch}{1.3}\begin{tabular}{|c|C{0.9cm}|C{0.9cm}|C{0.9cm}|C{0.9cm}|C{0.9cm}|C{0.9cm}|}\hline $x_i$ & $0$ & $1$ & $2$ & $3$ & $4$ & $5$\\\hline $F_i$ & $2$ & $7$ & $15$ & $22$ & $26$ & $28$\\\hline\end{tabular}}\end{center}

\ex{2} Hemos lanzado un dado 60 veces. Ha salido el 1 nueve veces, el 2 once, el 3 diez, el 4 ocho, el 5 doce y el 6 el resto. ¿Cuántas veces ha salido el 6? ¿Qué porcentaje de los lanzamientos ha salido un número par? ¿Y un número mayor que 4?

\ex{3} Una variable toma tres valores, $A$, $B$ y $C$. La frecuencia relativa de $A$ es $0{,}25$; la de $B$, $0{,}4$, y la frecuencia absoluta de $B$ es 16. Calcula $N$, la frecuencia relativa de $C$ y todas las frecuencias absolutas.

\ex{3} En una tabla de frecuencias, $F_3=12$, $F_4=19$ y $F_5=25$, y 5 es el valor mayor. ¿Cuántos datos hay? ¿Cuál es la frecuencia absoluta del 4? ¿Qué porcentaje de los datos es mayor que 3?


% =====================================================
\seccio{3}{Tablas con datos agrupados en intervalos}

\ex{1} Los intervalos de una tabla son $[10,\ 15)$, $[15,\ 20)$, $[20,\ 25)$ y $[25,\ 30)$. ¿En qué intervalo va cada uno de estos datos? ¿Cuál es la marca de clase de cada intervalo?\begin{tasks}(4)
  \task $15$
  \task $24{,}9$
  \task $25$
  \task $10$
\end{tasks}

\ex{1} Estos son los pesos (en kg) de 24 alumnos: \ $48,\ \allowbreak 55,\ \allowbreak 61,\ \allowbreak 44,\ \allowbreak 52,\ \allowbreak 58,\ \allowbreak 63,\ \allowbreak 47,\ \allowbreak 50,\ \allowbreak 56,\ \allowbreak 59,\ \allowbreak 66,\ \allowbreak 45,\ \allowbreak 53,\ \allowbreak 57,\ \allowbreak 60,\ \allowbreak 49,\ \allowbreak 54,\ \allowbreak 62,\ \allowbreak 51,\ \allowbreak 58,\ \allowbreak 64,\ \allowbreak 46,\ \allowbreak 55$. \ Agrúpalos en intervalos de amplitud 5 desde 40\,kg y haz la tabla con las marcas de clase, $f_i$, $F_i$ y porcentajes. ¿Qué porcentaje pesa menos de 55\,kg?\recorda{El intervalo $[a,\ b)$ contiene los datos $\geq a$ y $<b$. La \textbf{marca de clase} es el punto medio: $\frac{a+b}{2}$.}

\ex{2} Las puntuaciones de 40 alumnos en un test han sido: \ $52,\ \allowbreak 67,\ \allowbreak 45,\ \allowbreak 73,\ \allowbreak 38,\ \allowbreak 61,\ \allowbreak 55,\ \allowbreak 80,\ \allowbreak 49,\ \allowbreak 64,\ \allowbreak 58,\ \allowbreak 71,\ \allowbreak 43,\ \allowbreak 66,\ \allowbreak 59,\ \allowbreak 35,\ \allowbreak 77,\ \allowbreak 62,\ \allowbreak 50,\ \allowbreak 69,\ \allowbreak 57,\ \allowbreak 84,\ \allowbreak 47,\ \allowbreak 63,\ \allowbreak 54,\ \allowbreak 74,\ \allowbreak 41,\ \allowbreak 68,\ \allowbreak 60,\ \allowbreak 56,\ \allowbreak 31,\ \allowbreak 79,\ \allowbreak 65,\ \allowbreak 51,\ \allowbreak 88,\ \allowbreak 70,\ \allowbreak 46,\ \allowbreak 58,\ \allowbreak 62,\ \allowbreak 53$. \ Decide cuántos intervalos haces y de qué amplitud (usa la regla de la raíz) y haz la tabla de frecuencias. ¿Qué porcentaje ha sacado menos de 50 puntos?\recorda{Rango $=x_{\max}-x_{\min}$. \ Número de intervalos $\approx\sqrt N$. \ Amplitud $\approx\frac{\text{rango}}{\text{n.\textsuperscript{o} de intervalos}}$, redondeada hacia arriba.}

\ex{2} La tabla recoge las alturas (en cm) de 50 alumnos. Haz la columna de las $F_i$ y responde: ¿cuántos alumnos miden menos de 165\,cm? ¿Qué porcentaje mide 170\,cm o más? ¿Cuántos miden entre 155 y 170\,cm?\par\smallskip\begin{center}{\renewcommand{\arraystretch}{1.3}\begin{tabular}{|c|C{1.6cm}|C{1.6cm}|C{1.6cm}|C{1.6cm}|C{1.6cm}|C{1.6cm}|}\hline Altura & $[150,\ 155)$ & $[155,\ 160)$ & $[160,\ 165)$ & $[165,\ 170)$ & $[170,\ 175)$ & $[175,\ 180)$\\\hline $f_i$ & $4$ & $9$ & $15$ & $12$ & $7$ & $3$\\\hline\end{tabular}}\end{center}

\ex{3} Las marcas de clase de una tabla son 4, 8, 12 y 16, y todos los intervalos tienen la misma amplitud. ¿Cuáles son los intervalos? ¿Y si un intervalo tiene marca $32{,}5$ y amplitud 5?


% =====================================================
\seccio{4}{Gráficos estadísticos}

\ex{1} Dibuja el diagrama de barras y el polígono de frecuencias de los libros leídos por los 24 alumnos de la sección 2 (ejercicio 6).\recorda{En un \textbf{diagrama de barras}, las barras van separadas y tienen la altura de la frecuencia. El \textbf{polígono} une los extremos de arriba de las barras.}

\ex{1} Los 36 alumnos de una clase vienen al instituto así: a pie 15, en autobús 12, en coche 6 y en bici 3. Calcula el ángulo de cada sector y dibuja el diagrama de sectores.\recorda{Ángulo del sector $=\dfrac{f_i}{N}\cdot360^\circ$.}

\ex{2} Este diagrama muestra cuántas veces han viajado en avión los alumnos de una clase. Haz la tabla de frecuencias. ¿Cuántos alumnos hay? ¿Cuál es la moda? ¿Cuántos han viajado 3 o más veces?\par\smallskip\begin{center}\begin{tikzpicture}[baseline=(current bounding box.center),line cap=round,line join=round]
\draw[muted!30,line width=0.3pt] (0,0.000)--(4.700,0.000); \node[left,font=\tiny,text=muted,inner sep=1.5pt] at (0,0.000) {$0$};
\draw[muted!30,line width=0.3pt] (0,0.350)--(4.700,0.350); \node[left,font=\tiny,text=muted,inner sep=1.5pt] at (0,0.350) {$1$};
\draw[muted!30,line width=0.3pt] (0,0.700)--(4.700,0.700); \node[left,font=\tiny,text=muted,inner sep=1.5pt] at (0,0.700) {$2$};
\draw[muted!30,line width=0.3pt] (0,1.050)--(4.700,1.050); \node[left,font=\tiny,text=muted,inner sep=1.5pt] at (0,1.050) {$3$};
\draw[muted!30,line width=0.3pt] (0,1.400)--(4.700,1.400); \node[left,font=\tiny,text=muted,inner sep=1.5pt] at (0,1.400) {$4$};
\draw[muted!30,line width=0.3pt] (0,1.750)--(4.700,1.750); \node[left,font=\tiny,text=muted,inner sep=1.5pt] at (0,1.750) {$5$};
\draw[muted!30,line width=0.3pt] (0,2.100)--(4.700,2.100); \node[left,font=\tiny,text=muted,inner sep=1.5pt] at (0,2.100) {$6$};
\draw[muted!30,line width=0.3pt] (0,2.450)--(4.700,2.450); \node[left,font=\tiny,text=muted,inner sep=1.5pt] at (0,2.450) {$7$};
\draw[muted!30,line width=0.3pt] (0,2.800)--(4.700,2.800); \node[left,font=\tiny,text=muted,inner sep=1.5pt] at (0,2.800) {$8$};
\draw[-latex,muted,line width=0.6pt] (0,0)--(5.000,0) node[right,font=\scriptsize,text=navy]{Viajes};
\draw[-latex,muted,line width=0.6pt] (0,0)--(0,3.150) node[above,font=\scriptsize,text=navy]{Alumnos};
\draw[teal,fill=teal!35] (0.280,0) rectangle (0.820,1.750);
\node[below,font=\scriptsize,inner sep=2pt] at (0.550,0) {0};
\draw[teal,fill=teal!35] (1.180,0) rectangle (1.720,2.800);
\node[below,font=\scriptsize,inner sep=2pt] at (1.450,0) {1};
\draw[teal,fill=teal!35] (2.080,0) rectangle (2.620,2.100);
\node[below,font=\scriptsize,inner sep=2pt] at (2.350,0) {2};
\draw[teal,fill=teal!35] (2.980,0) rectangle (3.520,1.400);
\node[below,font=\scriptsize,inner sep=2pt] at (3.250,0) {3};
\draw[teal,fill=teal!35] (3.880,0) rectangle (4.420,0.700);
\node[below,font=\scriptsize,inner sep=2pt] at (4.150,0) {4};
\end{tikzpicture}\end{center}

\ex{2} Dibuja el histograma y el polígono de frecuencias de las alturas de los 50 alumnos de la sección 3 (ejercicio 16).\recorda{En un \textbf{histograma} las barras ocupan todo el intervalo y van juntas. El polígono une los puntos de arriba sobre las marcas de clase y se cierra en el eje.}

\ex{2} Se ha preguntado a 120 personas dónde prefieren pasar las vacaciones. ¿Cuántas personas han elegido cada lugar?\par\smallskip\begin{center}\begin{tikzpicture}[baseline=(current bounding box.center)]
\fill[teal!70,draw=white,line width=0.8pt] (0,0)--(90.000:1.45) arc[start angle=90.000,end angle=-18.000,radius=1.45cm]--cycle;
\node[font=\scriptsize\bfseries,text=white] at (36.000:0.899) {$108^\circ$};
\fill[accent!70,draw=white,line width=0.8pt] (0,0)--(-18.000:1.45) arc[start angle=-18.000,end angle=-108.000,radius=1.45cm]--cycle;
\node[font=\scriptsize\bfseries,text=white] at (-63.000:0.899) {$90^\circ$};
\fill[violet!70,draw=white,line width=0.8pt] (0,0)--(-108.000:1.45) arc[start angle=-108.000,end angle=-180.000,radius=1.45cm]--cycle;
\node[font=\scriptsize\bfseries,text=white] at (-144.000:0.899) {$72^\circ$};
\fill[gold!70,draw=white,line width=0.8pt] (0,0)--(-180.000:1.45) arc[start angle=-180.000,end angle=-240.000,radius=1.45cm]--cycle;
\node[font=\scriptsize\bfseries,text=white] at (-210.000:0.899) {$60^\circ$};
\fill[navy!70,draw=white,line width=0.8pt] (0,0)--(-240.000:1.45) arc[start angle=-240.000,end angle=-270.000,radius=1.45cm]--cycle;
\node[font=\scriptsize\bfseries,text=white] at (-255.000:0.899) {$30^\circ$};
\fill[teal!70] (1.85,0.590) rectangle ++(0.22,0.22); \node[right,font=\scriptsize] at (2.10,0.700) {Playa};
\fill[accent!70] (1.85,0.240) rectangle ++(0.22,0.22); \node[right,font=\scriptsize] at (2.10,0.350) {Montaña};
\fill[violet!70] (1.85,-0.110) rectangle ++(0.22,0.22); \node[right,font=\scriptsize] at (2.10,0.000) {Ciudad};
\fill[gold!70] (1.85,-0.460) rectangle ++(0.22,0.22); \node[right,font=\scriptsize] at (2.10,-0.350) {Pueblo};
\fill[navy!70] (1.85,-0.810) rectangle ++(0.22,0.22); \node[right,font=\scriptsize] at (2.10,-0.700) {Extranjero};
\end{tikzpicture}\end{center}

\ex{3} ¿Qué gráfico elegirías para representar cada variable? Justifícalo.\begin{tasks}(1)
  \task el número de hermanos de 30 alumnos;
  \task la altura de 200 personas;
  \task el partido que votan 500 personas;
  \task los minutos que tardan 50 alumnos en hacer los deberes.
\end{tasks}

\ex{3} Tres candidatos han sacado 96, 92 y 98 votos. Un cartel los representa con un diagrama de barras en el que el eje vertical empieza en 90. ¿Por qué es engañoso? Dibújalo bien.


% =====================================================
\seccio{5}{Parámetros de centralización}

\ex{1} Calcula la moda, la media y la mediana de: \ $7,\ \allowbreak 4,\ \allowbreak 9,\ \allowbreak 7,\ \allowbreak 3,\ \allowbreak 8,\ \allowbreak 7,\ \allowbreak 5,\ \allowbreak 6$.\recorda{\textbf{Moda}: el valor más repetido. \textbf{Media}: $\bar x=\dfrac{\text{suma de los datos}}{N}$. \textbf{Mediana}: el valor central de los datos \textbf{ordenados}; si hay dos, su media.}

\ex{1} Calcula la moda, la media y la mediana de: \ $12,\ \allowbreak 15,\ \allowbreak 9,\ \allowbreak 20,\ \allowbreak 14,\ \allowbreak 16,\ \allowbreak 11,\ \allowbreak 15$.

\ex{2} La tabla recoge cuántas veces han ido al cine en los últimos tres meses 30 alumnos. Calcula la moda, la media y la mediana.\par\smallskip\begin{center}{\renewcommand{\arraystretch}{1.3}\begin{tabular}{|c|C{0.9cm}|C{0.9cm}|C{0.9cm}|C{0.9cm}|C{0.9cm}|}\hline $x_i$ & $1$ & $2$ & $3$ & $4$ & $5$\\\hline $f_i$ & $4$ & $7$ & $9$ & $6$ & $4$\\\hline\end{tabular}}\end{center}\recorda{Con una tabla: $\bar x=\dfrac{\sum x_i\cdot f_i}{N}$. Para la mediana, mira en la columna $F_i$ dónde llega la posición central.}

\ex{2} Calcula la media aproximada de las alturas de los 50 alumnos de la sección 3 (ejercicio 16). ¿Cuál es el intervalo modal? ¿Y el intervalo mediano?

\ex{2} Pau ha sacado 6, 7 y 5 en tres exámenes. ¿Qué nota necesita en el cuarto para tener una media de $6{,}5$?\recorda{Si conocemos la media y el número de datos, conocemos la suma: \ $\text{suma}=\bar x\cdot N$.}

\ex{2} En una clase de 25 alumnos la nota media de un examen ha sido $6{,}4$ y en otra de 15 alumnos, $7{,}2$. ¿Cuál es la nota media de los 40 alumnos?

\ex{3} En una empresa pequeña, cuatro personas cobran 1\,100, 1\,200, 1\,250 y 1\,300\,€ al mes, y el gerente, 4\,650\,€. Calcula la media y la mediana de los sueldos. ¿Qué parámetro representa mejor el sueldo de un trabajador?

\ex{3} En esta tabla falta una frecuencia. Encuéntrala sabiendo que la media es $1{,}4$.\par\smallskip\begin{center}{\renewcommand{\arraystretch}{1.3}\begin{tabular}{|c|C{0.9cm}|C{0.9cm}|C{0.9cm}|C{0.9cm}|}\hline $x_i$ & $0$ & $1$ & $2$ & $3$\\\hline $f_i$ & $4$ & $a$ & $6$ & $2$\\\hline\end{tabular}}\end{center}


% =====================================================
\seccio{6}{Parámetros de dispersión}

\ex{1} Calcula el rango y la desviación media de: \ $3,\ \allowbreak 7,\ \allowbreak 8,\ \allowbreak 10,\ \allowbreak 12$.\recorda{Rango $=x_{\max}-x_{\min}$. \ $DM=\dfrac{\sum|x_i-\bar x|}{N}$.}

\ex{1} Calcula la varianza y la desviación típica de: \ $2,\ \allowbreak 4,\ \allowbreak 5,\ \allowbreak 7,\ \allowbreak 7$.\recorda{$\sigma^2=\dfrac{\sum(x_i-\bar x)^2}{N}$ \ o bien \ $\sigma^2=\dfrac{\sum x_i^2}{N}-\bar x^{\,2}$. \ $\sigma=\sqrt{\sigma^2}$.}

\ex{2} Calcula la media, la varianza y la desviación típica del número de hermanos de 20 alumnos:\par\smallskip\begin{center}{\renewcommand{\arraystretch}{1.3}\begin{tabular}{|c|C{0.9cm}|C{0.9cm}|C{0.9cm}|C{0.9cm}|C{0.9cm}|}\hline $x_i$ & $0$ & $1$ & $2$ & $3$ & $4$\\\hline $f_i$ & $5$ & $8$ & $4$ & $2$ & $1$\\\hline\end{tabular}}\end{center}

\ex{2} Dos jugadores de baloncesto han hecho estos puntos en cinco partidos. A: $6,\ 7,\ 7,\ 8,\ 7$; \ B: $10,\ 4,\ 9,\ 5,\ 7$. Calcula la media y la desviación típica de cada uno. ¿Quién es más regular?

\ex{2} La altura de un grupo de personas tiene $\bar x=160$\,cm y $\sigma=8$\,cm. El peso de un grupo de perros tiene $\bar x=40$\,kg y $\sigma=4$\,kg. ¿Qué grupo tiene más dispersión relativa?\recorda{Para comparar dispersiones con medias o unidades diferentes: $CV=\dfrac{\sigma}{\bar x}$ (en \%).}

\ex{3} Calcula la desviación típica y el coeficiente de variación de las alturas de los 50 alumnos (sección 3, ejercicio 16).

\ex{3} Los datos $2,\ 4,\ 5,\ 7,\ 7$ tienen $\bar x=5$ y $\sigma\approx1{,}90$. Sin volver a hacer todos los cálculos, encuentra la media, la desviación típica y el coeficiente de variación si:\begin{tasks}(2)
  \task sumamos 5 a todos los datos;
  \task multiplicamos todos los datos por 2.
\end{tasks}

\ex{3} Dos máquinas llenan botellas de un litro. La A llena de media 1\,002\,ml con $\sigma=3$\,ml y la B, 998\,ml con $\sigma=8$\,ml. ¿Qué máquina funciona mejor? Razónalo con el coeficiente de variación.


% =====================================================
\seccio{7}{Parámetros de posición: cuartiles y percentiles}

\ex{1} Calcula los cuartiles y el rango intercuartílico de: \ $4,\ \allowbreak 8,\ \allowbreak 15,\ \allowbreak 16,\ \allowbreak 23,\ \allowbreak 42,\ \allowbreak 7,\ \allowbreak 10,\ \allowbreak 12,\ \allowbreak 19,\ \allowbreak 30$.\recorda{Ordena los datos y calcula la posición $\frac{k\cdot N}{100}$ ($Q_1$: $\frac N4$; $Q_3$: $\frac{3N}{4}$). Si no es entera, coge el dato de la posición siguiente; si es entera, haz la media de ese dato y el siguiente.}

\ex{1} Explica qué quiere decir cada frase:\begin{tasks}(1)
  \task Un bebé está en el percentil 70 de peso.
  \task La nota de Lluís es el primer cuartil de la clase.
\end{tasks}

\ex{2} Las edades de los 12 miembros de un club son: \ $21,\ \allowbreak 25,\ \allowbreak 18,\ \allowbreak 30,\ \allowbreak 27,\ \allowbreak 22,\ \allowbreak 19,\ \allowbreak 26,\ \allowbreak 24,\ \allowbreak 28,\ \allowbreak 20,\ \allowbreak 23$. \ Calcula los cuartiles y los percentiles 10 y 90.

\ex{2} La tabla recoge los goles marcados por un equipo en 30 partidos. Calcula los cuartiles y los percentiles 10 y 90.\par\smallskip\begin{center}{\renewcommand{\arraystretch}{1.3}\begin{tabular}{|c|C{0.9cm}|C{0.9cm}|C{0.9cm}|C{0.9cm}|C{0.9cm}|C{0.9cm}|}\hline $x_i$ & $0$ & $1$ & $2$ & $3$ & $4$ & $5$\\\hline $f_i$ & $3$ & $8$ & $10$ & $6$ & $2$ & $1$\\\hline\end{tabular}}\end{center}\recorda{En una tabla, busca el primer valor con $F_i\geq\frac{k\cdot N}{100}$. Si una $F_i$ coincide exactamente, haz la media de ese valor y el siguiente.}

\ex{2} Con la tabla del ejercicio anterior, ¿a qué percentil corresponde, aproximadamente, marcar 2 goles? ¿Cuántos goles hay que marcar para estar en el cuartil 3 o por encima?

\ex{3} ¿En qué intervalo están el primer cuartil, el tercer cuartil y el percentil 90 de las alturas de los 50 alumnos de la sección 3 (ejercicio 16)?

\ex{3} En unas oposiciones con 250 aspirantes, solo aprueban quienes tienen una nota igual o superior al percentil 80. ¿Cuántos aspirantes aprueban, aproximadamente? Si Marta ha sacado una nota que es el percentil 75, ¿aprueba?


% =====================================================
\seccio{8}{Diagrama de cajas}

\ex{1} Dibuja el diagrama de cajas de los datos del ejercicio 41 (sección 7).\recorda{Cinco números: mínimo, $Q_1$, $Me$, $Q_3$ y máximo. La \textbf{caja} va de $Q_1$ a $Q_3$ con una línea en la mediana; los \textbf{bigotes} llegan al mínimo y al máximo.}

\ex{1} Este diagrama de cajas representa las edades de 80 personas que han ido a un concierto. Encuentra el rango y el rango intercuartílico. ¿Cuántas personas tienen más de 50 años? ¿Y entre 35 y 70 años?\par\smallskip\begin{center}\begin{tikzpicture}[baseline=(current bounding box.center),line cap=round]
\draw[muted!30,line width=0.3pt] (0.000,0)--(0.000,0.950); \draw[muted] (0.000,0.06)--(0.000,-0.06) node[below,font=\tiny,text=muted,inner sep=1.5pt]{$15$};
\draw[muted!30,line width=0.3pt] (0.667,0)--(0.667,0.950); \draw[muted] (0.667,0.06)--(0.667,-0.06) node[below,font=\tiny,text=muted,inner sep=1.5pt]{$20$};
\draw[muted!30,line width=0.3pt] (1.333,0)--(1.333,0.950); \draw[muted] (1.333,0.06)--(1.333,-0.06) node[below,font=\tiny,text=muted,inner sep=1.5pt]{$25$};
\draw[muted!30,line width=0.3pt] (2.000,0)--(2.000,0.950); \draw[muted] (2.000,0.06)--(2.000,-0.06) node[below,font=\tiny,text=muted,inner sep=1.5pt]{$30$};
\draw[muted!30,line width=0.3pt] (2.667,0)--(2.667,0.950); \draw[muted] (2.667,0.06)--(2.667,-0.06) node[below,font=\tiny,text=muted,inner sep=1.5pt]{$35$};
\draw[muted!30,line width=0.3pt] (3.333,0)--(3.333,0.950); \draw[muted] (3.333,0.06)--(3.333,-0.06) node[below,font=\tiny,text=muted,inner sep=1.5pt]{$40$};
\draw[muted!30,line width=0.3pt] (4.000,0)--(4.000,0.950); \draw[muted] (4.000,0.06)--(4.000,-0.06) node[below,font=\tiny,text=muted,inner sep=1.5pt]{$45$};
\draw[muted!30,line width=0.3pt] (4.667,0)--(4.667,0.950); \draw[muted] (4.667,0.06)--(4.667,-0.06) node[below,font=\tiny,text=muted,inner sep=1.5pt]{$50$};
\draw[muted!30,line width=0.3pt] (5.333,0)--(5.333,0.950); \draw[muted] (5.333,0.06)--(5.333,-0.06) node[below,font=\tiny,text=muted,inner sep=1.5pt]{$55$};
\draw[muted!30,line width=0.3pt] (6.000,0)--(6.000,0.950); \draw[muted] (6.000,0.06)--(6.000,-0.06) node[below,font=\tiny,text=muted,inner sep=1.5pt]{$60$};
\draw[muted!30,line width=0.3pt] (6.667,0)--(6.667,0.950); \draw[muted] (6.667,0.06)--(6.667,-0.06) node[below,font=\tiny,text=muted,inner sep=1.5pt]{$65$};
\draw[muted!30,line width=0.3pt] (7.333,0)--(7.333,0.950); \draw[muted] (7.333,0.06)--(7.333,-0.06) node[below,font=\tiny,text=muted,inner sep=1.5pt]{$70$};
\draw[muted!30,line width=0.3pt] (8.000,0)--(8.000,0.950); \draw[muted] (8.000,0.06)--(8.000,-0.06) node[below,font=\tiny,text=muted,inner sep=1.5pt]{$75$};
\draw[muted,line width=0.6pt] (0,0)--(8.000,0);
\node[right,font=\scriptsize,text=navy] at (8.100,0) {años};
\draw[thick] (0.667,0.550)--(2.667,0.550) (4.667,0.550)--(7.333,0.550);
\draw[thick] (0.667,0.420)--(0.667,0.680) (7.333,0.420)--(7.333,0.680);
\draw[teal,thick,fill=teal!25] (2.667,0.330) rectangle (4.667,0.770);
\draw[accent,line width=1.6pt] (4.000,0.330)--(4.000,0.770);
\end{tikzpicture}\end{center}

\ex{2} Las notas de un examen de dos clases se resumen en estos diagramas de cajas. Responde y razona:\par\smallskip\begin{center}\begin{tikzpicture}[baseline=(current bounding box.center),line cap=round]
\draw[muted!30,line width=0.3pt] (0.000,0)--(0.000,1.800); \draw[muted] (0.000,0.06)--(0.000,-0.06) node[below,font=\tiny,text=muted,inner sep=1.5pt]{$0$};
\draw[muted!30,line width=0.3pt] (0.800,0)--(0.800,1.800); \draw[muted] (0.800,0.06)--(0.800,-0.06) node[below,font=\tiny,text=muted,inner sep=1.5pt]{$1$};
\draw[muted!30,line width=0.3pt] (1.600,0)--(1.600,1.800); \draw[muted] (1.600,0.06)--(1.600,-0.06) node[below,font=\tiny,text=muted,inner sep=1.5pt]{$2$};
\draw[muted!30,line width=0.3pt] (2.400,0)--(2.400,1.800); \draw[muted] (2.400,0.06)--(2.400,-0.06) node[below,font=\tiny,text=muted,inner sep=1.5pt]{$3$};
\draw[muted!30,line width=0.3pt] (3.200,0)--(3.200,1.800); \draw[muted] (3.200,0.06)--(3.200,-0.06) node[below,font=\tiny,text=muted,inner sep=1.5pt]{$4$};
\draw[muted!30,line width=0.3pt] (4.000,0)--(4.000,1.800); \draw[muted] (4.000,0.06)--(4.000,-0.06) node[below,font=\tiny,text=muted,inner sep=1.5pt]{$5$};
\draw[muted!30,line width=0.3pt] (4.800,0)--(4.800,1.800); \draw[muted] (4.800,0.06)--(4.800,-0.06) node[below,font=\tiny,text=muted,inner sep=1.5pt]{$6$};
\draw[muted!30,line width=0.3pt] (5.600,0)--(5.600,1.800); \draw[muted] (5.600,0.06)--(5.600,-0.06) node[below,font=\tiny,text=muted,inner sep=1.5pt]{$7$};
\draw[muted!30,line width=0.3pt] (6.400,0)--(6.400,1.800); \draw[muted] (6.400,0.06)--(6.400,-0.06) node[below,font=\tiny,text=muted,inner sep=1.5pt]{$8$};
\draw[muted!30,line width=0.3pt] (7.200,0)--(7.200,1.800); \draw[muted] (7.200,0.06)--(7.200,-0.06) node[below,font=\tiny,text=muted,inner sep=1.5pt]{$9$};
\draw[muted!30,line width=0.3pt] (8.000,0)--(8.000,1.800); \draw[muted] (8.000,0.06)--(8.000,-0.06) node[below,font=\tiny,text=muted,inner sep=1.5pt]{$10$};
\draw[muted,line width=0.6pt] (0,0)--(8.000,0);
\node[right,font=\scriptsize,text=navy] at (8.100,0) {nota};
\draw[thick] (2.400,1.400)--(4.000,1.400) (5.600,1.400)--(7.200,1.400);
\draw[thick] (2.400,1.270)--(2.400,1.530) (7.200,1.270)--(7.200,1.530);
\draw[teal,thick,fill=teal!25] (4.000,1.180) rectangle (5.600,1.620);
\draw[accent,line width=1.6pt] (4.800,1.180)--(4.800,1.620);
\node[left,font=\small\bfseries] at (-0.15,1.400) {A};
\draw[thick] (0.800,0.550)--(3.200,0.550) (6.400,0.550)--(8.000,0.550);
\draw[thick] (0.800,0.420)--(0.800,0.680) (8.000,0.420)--(8.000,0.680);
\draw[accent,thick,fill=accent!25] (3.200,0.330) rectangle (6.400,0.770);
\draw[accent,line width=1.6pt] (4.800,0.330)--(4.800,0.770);
\node[left,font=\small\bfseries] at (-0.15,0.550) {B};
\end{tikzpicture}\end{center}\begin{tasks}(1)
  \task ¿Qué clase tiene la mediana más alta?
  \task ¿Qué clase tiene las notas más homogéneas?
  \task ¿En qué clase ha aprobado (5 o más) un porcentaje mayor de alumnos?
\end{tasks}

\ex{2} Dibuja el diagrama de cajas de los goles del ejercicio 44 (sección 7). ¿Qué indica que la media caja de la derecha y la de la izquierda sean iguales?

\ex{2} De unos datos sabemos que el mínimo es 12, $Q_1=18$, $Me=21$, $Q_3=30$ y el máximo, 34. ¿Qué porcentaje de los datos es:\begin{tasks}(4)
  \task menor que 21?
  \task entre 18 y 30?
  \task mayor que 18?
  \task entre 30 y 34?
\end{tasks}

\ex{3} Comprueba si hay valores atípicos en los datos $3,\ \allowbreak 21,\ \allowbreak 22,\ \allowbreak 24,\ \allowbreak 25,\ \allowbreak 26,\ \allowbreak 27,\ \allowbreak 28,\ \allowbreak 30$. Dibuja su diagrama de cajas y calcula su media y su mediana. ¿Qué parámetro representa mejor los datos?\recorda{Un dato es \textbf{atípico} si es menor que $Q_1-1{,}5\cdot RI$ o mayor que $Q_3+1{,}5\cdot RI$.}

\ex{3} Inventa ocho datos que tengan este resumen: mínimo 1, $Q_1=3$, $Me=5$, $Q_3=8$, máximo 12.

\vfill
\begin{center}\color{muted}\small\faCheckCircle[regular]\ \ ¡Ánimo! Comprueba tus resultados con las herramientas interactivas de los apuntes de la unidad.\end{center}

\end{document}
